\section*{Overview}

Disjoint Set Union is the data structure for connectivity under merges. It keeps track of which elements belong to the
same component while supporting fast ``find the leader'' and ``merge the two components'' operations.

In competitive programming, DSU shows up everywhere: Kruskal, offline connectivity, component counting on grids,
small-to-large style reasoning, and rollback-based dynamic connectivity.

\section*{When to Use It}

Use it when:

\begin{itemize}
  \item components only merge and never split,
  \item queries ask whether two elements are in the same component,
  \item you want to maintain a simple statistic per component such as size,
  \item the graph evolves by adding edges rather than deleting them.
\end{itemize}

If the problem needs deletions or cut operations, plain DSU is usually the wrong tool and rollback DSU or a dynamic
tree structure becomes more relevant.

\section*{Core Idea}

Represent each component as a rooted tree. Every node points to a parent, and the root of each tree is the component
leader.

The two operations are:

\begin{itemize}
  \item \textbf{find(x)}: climb to the root and return the leader,
  \item \textbf{unite(a, b)}: connect the roots if they are different.
\end{itemize}

The tree itself is not the graph component. It is only an internal representation of the partition.

\section*{Key Insight}

The reason DSU is fast is the combination of two heuristics:

\begin{itemize}
  \item \textbf{path compression}: after finding the root, rewrite the visited path so future finds are shorter,
  \item \textbf{union by size or rank}: attach the smaller tree under the larger one.
\end{itemize}

Path compression makes repeated finds on the same region cheaper. Union by size keeps the internal trees from growing
deep in the first place. Together they make the amortized cost essentially constant for contest purposes.

\section*{Operations / Main Technique}

The standard interface is:

\begin{itemize}
  \item \textbf{find(x)},
  \item \textbf{unite(a, b)},
  \item \textbf{same(a, b)},
  \item \textbf{size(x)}.
\end{itemize}

\paragraph{Component invariant.}
Only the root stores component-wide information such as size. That is a useful mental rule: if a value should mean
``for the whole set'', keep it on the leader.

\paragraph{Worked example.}
If \(\{1, 2, 3\}\) and \(\{4, 5\}\) are separate components, uniting \(2\) and \(5\) does not attach those exact
vertices conceptually. It finds their leaders first, then merges the two whole components through those leaders.

\section*{Correctness Intuition}

At all times, every element belongs to exactly one leader tree. Two elements are in the same component exactly when
their leaders are equal.

Path compression does not change the partition, because it only shortcuts parent pointers inside one component toward
the same root. Union by size also preserves correctness, because it merges two different roots into one new component
without losing any membership relation.

\section*{Complexity Analysis}

With path compression and union by size/rank:

\begin{itemize}
  \item \texttt{find}, \texttt{unite}, \texttt{same}, \texttt{size}: \(O(\alpha(n))\) amortized,
  \item memory: \(O(n)\).
\end{itemize}

The inverse Ackermann function \(\alpha(n)\) grows so slowly that for contest input sizes it behaves like a tiny
constant.

\section*{Implementation}

The reference code stores one array \texttt{parent} where roots keep negative component sizes and non-roots keep their
parent index. That layout is compact and keeps all the basic operations together.

It also tracks the number of connected components explicitly, which is often convenient in graph problems where the
final answer depends on how many merges succeeded.

\section*{Common Pitfalls}

\begin{itemize}
  \item Forgetting to find the leaders before merging.
  \item Mixing ``rank'' and ``size'' logic halfway through one implementation.
  \item Updating component data on a non-root instead of the root.
  \item Assuming DSU can handle deletions just because connectivity is involved.
  \item Adding path compression to a rollback DSU, where that optimization becomes inconvenient or incorrect.
\end{itemize}

\section*{Variants / Extensions}

\begin{itemize}
  \item DSU rollback for offline dynamic connectivity,
  \item DSU with parity or xor distance information,
  \item DSU that stores custom data per component,
  \item offline painting / next-unprocessed-cell DSU tricks,
  \item union-by-size ``small to large'' ideas on explicit containers.
\end{itemize}

Rollback DSU is the extension I see most often in harder contest problems, especially when edges are added and removed
over time but the whole query sequence is known offline.

\section*{Practice Problems}

\begin{itemize}
  \item Kruskal minimum spanning tree.
  \item Grid component counting after activating cells.
  \item Offline connectivity where edges are processed in reverse time.
  \item Problems that require storing one extra statistic such as component sum or bipartite validity.
\end{itemize}

\section*{References}

\begin{itemize}
  \item \href{https://cp-algorithms.com/data_structures/disjoint_set_union.html}{cp-algorithms: Disjoint Set Union}.
  \item \href{https://atcoder.github.io/ac-library/production/document_en/dsu.html}{AtCoder Library: dsu}.
  \item \href{https://usaco.guide/CPH.pdf}{Competitive Programmer's Handbook}.
  \item \href{https://codeforces.com/catalog}{Codeforces Catalog}.
\end{itemize}
